\scanpage{83}
Soit $\deleted{\unsure{\widetilde{\mathcal T}\longrightarrow\mathcal T}}$
\[(1)\qquad\widetilde{\mathcal T}\longrightarrow\mathcal T\]
un morphisme de groupoïdes, et
\[(2)\qquad(\widetilde u_i)_{i\in I}\]
une famille de flèches de $\widetilde{\mathcal T}$
\[(3)\qquad(u_i)_{i\in I}\]
la famille correspondante de $\mathcal T$. On suppose que celle-ci est génératrice (donc $\operatorname{Ob}\mathcal T$ est formé des origines et ext. des flèches $u_i$) et on se donne une famille de relations \unsure{définissantes} de $\mathcal T$, avec les flèches génératrices $u_i$. Cela se fait, techniquement, en \deleted{\unsure{introduisant}}\inserted{\unsure{introduisant}} l'ens d'indices
\[(4)\qquad I'=I\times\{\pm1\}\]
et en définissant, pour $i'=(i,\varepsilon)\in I'$
\[(5)\qquad u'_{i'}=u_i^\varepsilon\]
et en se donnant une famille, \deleted{de} indexée par un ens d'indices $A$, d'éléments de l'ensemble somme
\[(6)\qquad\widehat I'=\coprod_{n\in\mathbb N^*}I'^n.\]
Donc pour $\forall\alpha\in A$, on a un $n_\alpha\in\mathbb N^*$,
\mspage{82}
